
\subsection*{4.5}
On
appelle équivalence numérique la relation d'équivalence définie par le produit d'intersection @@M0@@ :

\textbf{Définition 4.5.1.} \uline{On dit que deux éléments} @@M1@@ \uline{et} @@M2@@ \uline{de} @@M3@@ @@M4@@ (resp. \uline{de} @@M5@@\uline{) sont numériquement équivalents, s'ils définissent le même homomorphisme} @@M6@@ (resp. @@M7@@\uline{).
On dit que deux} @@M8@@\uline{-modules inversibles} @@M9@@ et @@M10@@ \uline{sont numériquement équivalents, si leurs premières classes de Chern} @@M11@@ \uline{sont numériquement équivalentes.}

\textbf{Proposition 4.5.2.} \uline{Un} @@M12@@\uline{-module inversible} @@M13@@ \uline{est numériquement équivalent à zéro si et seulement si}
@@M14@@
\uline{pour tout} @@M15@@ \uline{et pour toute courbe intègre fermée} @@M16@@.

Cela résulte de 1.1.2 et 4.3.1.

On désigne par @@M17@@ le sous-groupe de @@M18@@ composé des classes de @@M19@@-modules inversibles numériquement équivalents à zéro. De 2.4.4 on tire

\textbf{Corollaire 4.5.3.} @@M20@@.

On va démontrer dans Exp XIII que l'on a même @@M21@@.

\section*{5. L'isomorphisme @@M22@@}

\subsection*{5.1.}
Soit @@M23@@ un topos annelé quelconque. Alors on peut définir l'anneau @@M24@@ et la filtration @@M25@@ exactement comme pour les schémas. De même on définit le groupe @@M26@@ comme le groupe des classes d'isomorphie des @@M27@@-modules inversibles. On désigne par
